\documentclass[11pt,a4paper]{article}

\usepackage{amsmath,amssymb}
\usepackage{physics}
\usepackage{braket}
\usepackage{hyperref}
\usepackage[margin=2.5cm]{geometry}
\usepackage{graphicx}
\usepackage{slashed}

\hypersetup{
  colorlinks=true,
  linkcolor=blue,
  citecolor=blue,
  urlcolor=blue,
}

% ----------------------------------------------------------------------
% Custom notation
% ----------------------------------------------------------------------
\newcommand{\eps}{\varepsilon}
\newcommand{\Gammaint}{\Gamma_{\text{interp}}}
\newcommand{\Gammasink}{\Gamma_{\text{sink}}}
\newcommand{\Cg}{C\gamma_5}
\newcommand{\CgT}{C\gamma_t\gamma_5}
\newcommand{\CgZ}{C\gamma_z\gamma_5}
\newcommand{\Cmat}{C}
\newcommand{\xvec}{x}
\newcommand{\xzerovec}{x_0}
\newcommand{\bperp}{b_\perp}
\newcommand{\bvec}{b}
\newcommand{\ehat}[1]{\hat{e}_{#1}}
\newcommand{\Oop}{\mathcal{O}}

\title{Proton qTMD and PDF Measurement \\
       \large Physics and Contraction Note}
\author{PyQUDA Measurement Collaboration}
\date{\today}

\begin{document}

\maketitle
\tableofcontents

% ======================================================================
\section{Physics Overview}
% ======================================================================

Proton transverse momentum dependent (TMD) distributions encode the three-dimensional
momentum structure of the nucleon in QCD. Unlike the pion, the proton is a baryon
composed of three valence quarks, which introduces two key complications relative to
meson TMD measurements:

\begin{enumerate}
  \item The proton interpolating field requires a color-antisymmetric contraction via
    the Levi-Civita tensor $\eps_{abc}$, producing a baryon two-point function with
    two distinct contraction terms (direct and exchange).
  \item Sequential sources are flavor-dependent: the up-quark insertion involves four
    epsilon-contracted terms, while the down-quark insertion involves two.
\end{enumerate}

The leading-twist TMDs accessible from this measurement include:
\begin{itemize}
  \item $f_1(x, k_\perp^2)$ -- unpolarized quark distribution,
  \item $g_{1L}(x, k_\perp^2)$ -- helicity distribution,
  \item $h_1(x, k_\perp^2)$ -- transversity distribution,
\end{itemize}
where $x$ is the longitudinal momentum fraction and $k_\perp$ is the intrinsic
transverse momentum. The Fourier-conjugate of the TMD in impact-parameter space
$b_\perp$ is the quasi-TMD (qTMD) wave function, which is the primary object
computed in this lattice QCD measurement.

This document describes in full technical detail: (i) the Wick contraction of the
proton two-point function, (ii) the sequential-source construction for the connected
three-point qTMD function, (iii) the Wilson-line index generation for CG and
GI-displaced propagators, (iv) the einsum decomposition strategy used in the PyQUDA
GPU implementation, and (v) the HDF5 output conventions. The target audience is
lattice QCD practitioners who need to understand every contraction index.

% ======================================================================
\section{Proton Interpolating Field}
% ======================================================================

The proton interpolating operator used throughout this measurement is the standard
color-singlet three-quark field (suppressing the time argument for brevity):

\begin{equation}
  \chi_\alpha(x) =
    \eps_{abc}\,
    \bigl[ u_a^T(x) \; \Cmat \Gammaint \; d_b(x) \bigr] \;
    u_{c,\alpha}(x),
  \label{eq:proton_interp}
\end{equation}

where:
\begin{itemize}
  \item $\alpha \in \{0,1,2,3\}$ is the free Dirac index of the proton,
  \item $a,b,c \in \{0,1,2\}$ are fundamental color indices,
  \item $\eps_{abc}$ is the totally antisymmetric Levi-Civita tensor with
    $\eps_{012} = +1$,
  \item $\Cmat = i\gamma_2\gamma_4$ is the charge-conjugation matrix in the Euclidean
    Dirac basis,
  \item $\Gammaint$ selects the spin structure of the diquark.
\end{itemize}

Three interpolator choices are supported, denoted by string tags:

\begin{center}
  \begin{tabular}{c|c|c}
    Tag & Matrix & Expression \\
    \hline
    \texttt{"5"}  & $\Cg$   & $\Cmat \gamma_5$ \\
    \texttt{"T5"} & $\CgT$  & $\Cmat \gamma_t \gamma_5$ \\
    \texttt{"Z5"} & $\CgZ$  & $\Cmat \gamma_z \gamma_5$
  \end{tabular}
\end{center}

In the PyQUDA code (using the QUDA bitmask gamma basis), these are constructed as:
\begin{align}
  \Cg   &= (i\gamma_2 \gamma_8) \cdot \gamma_{15}, \label{eq:Cg5_code} \\
  \CgT  &= (i\gamma_2 \gamma_8) \cdot \gamma_7,  \label{eq:CgT5_code} \\
  \CgZ  &= (i\gamma_2 \gamma_8) \cdot \gamma_{11}, \label{eq:CgZ5_code}
\end{align}
where the QUDA basis uses $\gamma_8 = \gamma_t$, $\gamma_7 = \gamma_t\gamma_5$,
$\gamma_{15} = \gamma_5$, and $\gamma_{11} = \gamma_z\gamma_5$. The factor
$i\gamma_2\gamma_8$ is the charge conjugation matrix in this basis.

The conjugate field is:
\begin{equation}
  \bar{\chi}_\alpha(x) =
    \eps_{def}\,
    \bar{u}_{d,\alpha}(x) \;
    \bigl[ \bar{d}_e(x) \; \Gammaint^\dagger \Cmat^\dagger \;
           \bar{u}_f^T(x) \bigr],
  \label{eq:proton_bar}
\end{equation}
where $\Gammaint^\dagger = \gamma_4 \Gammaint^\dagger \gamma_4$ in Euclidean space
and $\Cmat^\dagger = \Cmat^{-1} = -\Cmat$.

% ======================================================================
\section{Proton Two-Point Function}
\label{sec:two_point}
% ======================================================================

\subsection{Definition}

The momentum-projected proton two-point correlation function with a sink spin
structure $\Gammasink$ is:

\begin{equation}
  C_2^{\Gammasink}(p, t) =
    \sum_{\xvec} \exp\bigl(-i p\cdot (\xvec - \xzerovec)\bigr) \;
    \mathcal{P}_{\alpha\alpha'} \;
    \bigl\langle
      \chi_\alpha(\xvec, t) \; \Gammasink \;
      \bar{\chi}_{\alpha'}(\xzerovec, 0)
    \bigr\rangle,
  \label{eq:C2_def}
\end{equation}

where $\mathcal{P}_{\alpha\alpha'}$ is the spin projection matrix (e.g.\ unpolarized
$P_+ = \frac{1}{4}(\gamma_0 + I)$ or polarized variants). The sum over $\xvec$ runs
over the full spatial lattice at fixed time slice $t$.

\subsection{Wick Contraction -- Complete Derivation}

Inserting the interpolating fields from Eqs.~\eqref{eq:proton_interp}
and~\eqref{eq:proton_bar} into the expectation value, and expanding the quark fields:

\begin{align}
  \bigl\langle \chi_\alpha(x) \; {\Gammasink}_{\alpha\alpha'} \;
    \bar{\chi}_{\alpha'}(x_0) \bigr\rangle
  &=
  \eps_{abc}\, \eps_{def} \;
  \bigl(\Cmat \Gammaint\bigr)_{ij}\,
  \bigl(\Cmat \Gammaint\bigr)_{kl}\,
  {\Gammasink}^{\alpha\alpha'}_{mn} \notag \\[4pt]
  &\quad\times
  \bigl\langle
    u_a^i(x)\, d_b^j(x)\, u_c^\alpha(x) \;
    \bar{u}_d^m(x_0)\, \bar{d}_e^k(x_0)\, \bar{u}_f^n(x_0)
  \bigr\rangle,
  \label{eq:raw_contraction}
\end{align}

where the spinor indices are denoted by $i,j,k,l,m,n \in \{0,1,2,3\}$ for the quark
bilinears, with $\alpha,\alpha'$ being the external proton spin indices. The color
indices are $a,\dots,f \in \{0,1,2\}$.

\medskip\noindent
\textbf{Step 1: Identify the quark flavors.}
The interpolating field contains two up quarks and one down quark:
\begin{itemize}
  \item $u_a$, $u_c$ are up quarks (labeled by their Dirac indices and color positions),
  \item $d_b$ is a down quark.
\end{itemize}
Similarly, the conjugate field contains $\bar{u}_d$, $\bar{u}_f$ (up) and $\bar{d}_e$
(down).

\medskip\noindent
\textbf{Step 2: Enumerate the Wick pairings.}
Wick's theorem instructs us to sum over all permutations of pairing each quark field
with each anti-quark field. Because the up quarks $u_a$ and $u_c$ are identical
fermions, there are two inequivalent ways to contract the up-quark operators:

\begin{center}
  \begin{tabular}{c|c|l}
    Term & Pairing & Description \\
    \hline
    Term~1 & $u_a \leftrightarrow \bar{u}_d$, $u_c \leftrightarrow \bar{u}_f$ & Direct \\
    Term~2 & $u_a \leftrightarrow \bar{u}_f$, $u_c \leftrightarrow \bar{u}_d$ & Exchange \\
  \end{tabular}
\end{center}

The down quark always contracts as $d_b \leftrightarrow \bar{d}_e$, since there is
exactly one down quark and one anti-down quark.

\medskip\noindent
\textbf{Step 3: Write each term with quark propagators.}
The quark propagator is defined as:
\begin{equation}
  S_q(x, x_0)^{i\alpha}_{ab} \equiv \langle q_a^i(x) \, \bar{q}_b^\alpha(x_0) \rangle,
\end{equation}
where $i,\alpha$ are spin indices and $a,b$ are color indices. For isospin-symmetric
computations ($m_u = m_d$), $S_u = S_d \equiv S$, so the same propagator is used
for all three valence quarks.

\medskip\noindent
\textbf{Term 1 (Direct):}
\begin{align}
  C_2^{(1)} &=
  \eps_{abc}\, \eps_{def}\,
  (\Cmat\Gammaint)_{ij}\,
  (\Cmat\Gammaint)_{kl}\,
  (\Gammasink)_{mn} \notag \\[4pt]
  &\quad\times
  S(x,x_0)^{ik}_{ad}\,
  S(x,x_0)^{jl}_{be}\,
  S(x,x_0)^{mn}_{cf},
  \label{eq:term1_full}
\end{align}
where the index assignments are:
\begin{itemize}
  \item $u_a$ (spin $i$, color $a$) $\to$ $\bar{u}_d$ (spin $k$, color $d$):
    propagator $S(x,x_0)^{ik}_{ad}$,
  \item $d_b$ (spin $j$, color $b$) $\to$ $\bar{d}_e$ (spin $l$, color $e$):
    propagator $S(x,x_0)^{jl}_{be}$,
  \item $u_c$ (spin $\alpha \equiv m$, color $c$) $\to$ $\bar{u}_f$ (spin $n$, color $f$):
    propagator $S(x,x_0)^{mn}_{cf}$.
\end{itemize}

\medskip\noindent
\textbf{Term 2 (Exchange):}
The exchange term swaps the two up-quark assignments:
\begin{align}
  C_2^{(2)} &=
  \eps_{abc}\, \eps_{def}\,
  (\Cmat\Gammaint)_{ij}\,
  (\Cmat\Gammaint)_{kl}\,
  (\Gammasink)_{mn} \notag \\[4pt]
  &\quad\times
  S(x,x_0)^{ik}_{ad}\,
  S(x,x_0)^{jn}_{be}\,
  S(x,x_0)^{ml}_{cf},
  \label{eq:term2_full}
\end{align}
where now:
\begin{itemize}
  \item $u_a$ (spin $i$, color $a$) $\to$ $\bar{u}_d$ (spin $k$, color $d$):
    same as Term~1,
  \item $d_b$ (spin $j$, color $b$) $\to$ $\bar{u}_f$ (spin $n$, color $f$):
    \emph{exchanged} -- $d$ contracts to $u$,
  \item $u_c$ (spin $m$, color $c$) $\to$ $\bar{d}_e$ (spin $l$, color $e$):
    \emph{exchanged} -- $u$ contracts to $d$.
\end{itemize}

Note that in the exchange term, the flavor labels $u$ and $d$ are interchanged on the
propagators (a down quark connects to the final $\bar{u}_f$, and an up quark connects
to the final $\bar{d}_e$). Under isospin symmetry, all propagators are identical, so the
flavor labels serve only to track the index routing.

\medskip\noindent
\textbf{Step 4: The sign convention.}
The fully antisymmetric contraction of two $\eps$ tensors with Fermion-exchange
antisymmetry produces an overall minus sign:
\begin{equation}
  C_2^{\Gammasink}(p, t) =
    -\Bigl( C_2^{(1)} + C_2^{(2)} \Bigr) \times
    \sum_{\xvec} e^{-ip\cdot(\xvec-\xzerovec)} \; \mathcal{P}_{\alpha\alpha'},
  \label{eq:C2_sign}
\end{equation}
where $C_2^{(1)}$ and $C_2^{(2)}$ are defined by Eqs.~\eqref{eq:term1_full}
and~\eqref{eq:term2_full} respectively. The negative sign originates from the fermion
anticommutation required to bring the quark operators into canonical contraction
order. After the Wick expansion, each term acquires a sign $(-1)^{n_P}$ where $n_P$
is the parity of the permutation mapping the original operator ordering to the
contracted pairing. Since each contraction is a product of three Wick pairs, the
permutation that aligns the color-epsilon structure with the contracted propagator
pairs yields $n_P = 1$ for both terms.

\medskip\noindent
\textbf{Step 5: Compact tensor notation.}
The two-point function can be written compactly in a factorized diagrammatic form.
Omitting the external momentum phase and spin projection for clarity, and using
position-space shorthand:

\begin{align}
  \text{Term~1:}&\quad
    \eps_{abc}\,\eps_{def}\,
    (\Cmat\Gammaint)_{ij}\,(\Cmat\Gammaint)_{kl}\,
    (\Gammasink)_{mn}\,
    S^{ik}_{ad}\,S^{jl}_{be}\,S^{mn}_{cf}, \\[4pt]
  \text{Term~2:}&\quad
    \eps_{abc}\,\eps_{def}\,
    (\Cmat\Gammaint)_{ij}\,(\Cmat\Gammaint)_{kl}\,
    (\Gammasink)_{mn}\,
    S^{ik}_{ad}\,S^{jn}_{be}\,S^{ml}_{cf}.
\end{align}

The key difference between the two terms is purely in the index routing of the third
propagator and the exchange of indices $(l,n)$ on the second and third propagators.

\subsection{Diagrammatic Interpretation}

Term~1 corresponds to the ``direct'' contraction:
\begin{center}
  \texttt{u\_a --[C$\Gamma$]-- d\_b} \quad \texttt{u\_c --[$\Gammasink$]-- ubar\_f} \\
  with $\eps_{abc}$ at source and $\eps_{def}$ at sink.
\end{center}

Term~2 corresponds to the ``exchange'' contraction where the two identical up quarks
are crossed:
\begin{center}
  \texttt{u\_a --[C$\Gamma$]-- d\_b} \quad (u and d sink legs exchanged).
\end{center}

There are exactly two terms because the up quarks are identical fermions but are
distinguished by their role in the interpolating field: one forms the scalar diquark
with the down quark (via $\Cmat\Gammaint$), the other carries the open spin index.
The two contractions differ in which up quark at the source connects to which up
anti-quark at the sink.

% ======================================================================
\section{Optimized Contraction Code}
\label{sec:contraction_code}
% ======================================================================

The function \texttt{contract\_2pt\_TMD} in \texttt{proton\_qTMD\_pyquda.py}
implements the two-point contraction on GPU accelerators. The full 8-index
tensor contraction (over $a,b,c,d,e,f,i,j,k,l,m,n$) is too large for a single
einsum call (the intermediate tensors would exceed GPU memory). The code therefore
splits the contraction into a sequence of smaller 2-tensor and 3-tensor einsum
operations, carefully managing intermediate memory through explicit deletion.

The propagator data is shaped as $[\omega, t, z, y, x, \mathrm{spin}_a,
\mathrm{spin}_b, \mathrm{color}_a, \mathrm{color}_b]$ where $\omega$ is the even-odd
parity index. The momentum phases array is indexed as $[p, \omega, t, z, y, x]$.
The 16 gamma matrices for the sink insertion form an array $[g, t_4, m, n]$ where
$g \in \{0,\dots,15\}$ labels the bilinear.

\subsection{Term 1 Decomposition}

The full Term~1 einsum would be:
\begin{equation}
  \eps_{abc}\,\eps_{def}\,
  (\Cmat\Gammaint)_{ij}\,(\Cmat\Gammaint)_{kl}\,
  (P_{2\mathrm{pt}})_{g t_4 mn}\,
  S^{ik}_{ad}\; S^{jl}_{be}\; S^{mn}_{cf}\;
  \Phi_p
  \;\to\; [g, p, t],
\end{equation}
where $\Phi_p$ is the momentum phase and $P_{2\mathrm{pt}}$ holds the 16 sink gammas.

The decomposition proceeds in three stages:

\medskip\noindent
\textbf{Stage 1 --- Sink block (split).}
The sink block contracts $\eps_{abc}$ with the first propagator and
$(\Cmat\Gammaint)_{ij}$ with the second propagator, then combines on indices
$i$ and $b$:
\begin{align}
  \texttt{t1\_s1} &: \quad \eps_{abc} \times S^{ik}_{ad}
    \;\to\; [\omega,t,z,y,x,i,k,b,c,d], \\
  \texttt{t1\_s2} &: \quad (\Cmat\Gammaint)_{ij} \times S^{jl}_{be}
    \;\to\; [\omega,t,z,y,x,i,l,b,e], \\
  \texttt{term1\_sink} &: \quad
    \texttt{t1\_s1} \times \texttt{t1\_s2} \;\to\;
    [\omega,t,z,y,x,k,l,c,d,e] \quad (\text{contract } i,b).
\end{align}
After this stage, the tensors \texttt{t1\_s1} and \texttt{t1\_s2} are deleted.

\medskip\noindent
\textbf{Stage 2 --- P3 block.}
The sink gamma array is contracted with the third propagator:
\begin{equation}
  \texttt{term1\_p3}: \quad
    (P_{2\mathrm{pt}})_{gmn} \times S^{mn}_{cf}
    \;\to\; [g,\omega,t,z,y,x,c,f].
\end{equation}

\medskip\noindent
\textbf{Stage 3 --- Final assembly (split).}
The sink block and P3 block are combined and multiplied by the momentum phases:
\begin{align}
  \texttt{t1\_f1} &: \quad
    \eps_{def} \times \texttt{term1\_sink}
    \;\to\; [\omega,t,z,y,x,k,l,c,f] \quad (\text{contract } d,e), \\
  \texttt{t1\_f2} &: \quad
    (\Cmat\Gammaint)_{kl} \times \texttt{t1\_f1}
    \;\to\; [\omega,t,z,y,x,c,f] \quad (\text{contract } k,l), \\
  \texttt{t1\_f3} &: \quad
    \texttt{t1\_f2} \times \texttt{term1\_p3}
    \;\to\; [g,\omega,t,z,y,x] \quad (\text{contract } c,f), \\
  \texttt{term1} &: \quad
    \Phi_p \times \texttt{t1\_f3}
    \;\to\; [g,p,t].
\end{align}
Each intermediate tensor is deleted immediately after use.

\subsection{Term 2 Decomposition}

Term~2 has a different index routing ($S^{jn}_{be}$ instead of $S^{jl}_{be}$ and
$S^{ml}_{cf}$ instead of $S^{mn}_{cf}$), so the decomposition differs:

\medskip\noindent
\textbf{Stage 1 --- Sink block.}
\begin{align}
  \texttt{t2\_s1} &: \quad \eps_{abc} \times S^{ik}_{ad}
    \;\to\; [\omega,t,z,y,x,i,k,b,c,d], \\
  \texttt{t2\_s2} &: \quad (\Cmat\Gammaint)_{ij} \times S^{jn}_{be}
    \;\to\; [\omega,t,z,y,x,i,n,b,e] \quad (\text{note: index } n \text{ not } l), \\
  \texttt{term2\_sink} &: \quad
    \texttt{t2\_s1} \times \texttt{t2\_s2}
    \;\to\; [\omega,t,z,y,x,k,n,c,d,e] \quad (\text{contract } i,b).
\end{align}

\medskip\noindent
\textbf{Stage 2 --- P3 block.}
The sink gamma contracts with the third propagator, keeping the exchange index $l$:
\begin{equation}
  \texttt{term2\_p3}: \quad
    (P_{2\mathrm{pt}})_{gmn} \times S^{ml}_{cf}
    \;\to\; [g,\omega,t,z,y,x,n,l,c,f].
\end{equation}
Note that $m$ is contracted between $P_{2\mathrm{pt}}$ and the propagator, while $n$
and $l$ remain open for the exchange contraction.

\medskip\noindent
\textbf{Stage 3 --- Final assembly.}
\begin{align}
  \texttt{t2\_f1} &: \quad
    \eps_{def} \times \texttt{term2\_sink}
    \;\to\; [\omega,t,z,y,x,k,n,c,f] \quad (\text{contract } d,e), \\
  \texttt{t2\_f2} &: \quad
    \texttt{t2\_f1} \times \texttt{term2\_p3}
    \;\to\; [g,\omega,t,z,y,x,k,l] \quad (\text{contract } n,c,f), \\
  \texttt{t2\_f3} &: \quad
    (\Cmat\Gammaint)_{kl} \times \texttt{t2\_f2}
    \;\to\; [g,\omega,t,z,y,x] \quad (\text{contract } k,l), \\
  \texttt{term2} &: \quad
    \Phi_p \times \texttt{t2\_f3}
    \;\to\; [g,p,t].
\end{align}

\subsection{Final Summation}

The complete two-point function is:
\begin{equation}
  \boxed{C_2[g, p, t] = -\bigl(\texttt{term1} + \texttt{term2}\bigr).}
\end{equation}

After the contraction, the result is MPI-gathered and saved as HDF5 (rank 0 only).
The saved data is time-rolled by the negative of the source time to align $t=0$
at the source.

% ======================================================================
\section{Gamma Matrix Conventions}
% ======================================================================

\subsection{Sink Bilinears}

Sixteen bilinear gamma insertions are scanned at the sink, ordered as:

\begin{center}
  \begin{tabular}{c|l}
    Index & Gamma \\
    \hline
    0  & 5   ($\gamma_5$) \\
    1  & T   ($\gamma_t$) \\
    2  & T5  ($\gamma_t\gamma_5$) \\
    3  & X   ($\gamma_x$) \\
    4  & X5  ($\gamma_x\gamma_5$) \\
    5  & Y   ($\gamma_y$) \\
    6  & Y5  ($\gamma_y\gamma_5$) \\
    7  & Z   ($\gamma_z$) \\
    8  & Z5  ($\gamma_z\gamma_5$) \\
    9  & I   (identity) \\
    10 & SXT ($\sigma_{xt}$) \\
    11 & SXY ($\sigma_{xy}$) \\
    12 & SXZ ($\sigma_{xz}$) \\
    13 & SYT ($\sigma_{yt}$) \\
    14 & SYZ ($\sigma_{yz}$) \\
    15 & SZT ($\sigma_{zt}$)
  \end{tabular}
\end{center}

In the QUDA bitmask gamma basis, these correspond to the gamma matrix ID array:
\begin{equation}
  \texttt{pyquda\_gammas\_order} =
    [15, 8, 7, 1, 14, 2, 13, 4, 11, 0, 9, 3, 5, 10, 6, 12].
\end{equation}

\subsection{Interpolator Matrices}

The proton interpolator matrices $C\Gammaint$ described in Section~2 are
pre-computed and stored as $4 \times 4$ complex arrays. In the QUDA basis:

\begin{itemize}
  \item \texttt{Cg5}  = $(i\gamma_2 \gamma_8) \cdot \gamma_{15}$
    $\;\leftrightarrow\;$ $C\gamma_5$,
  \item \texttt{CgT5} = $(i\gamma_2 \gamma_8) \cdot \gamma_7$
    $\;\leftrightarrow\;$ $C\gamma_t\gamma_5$,
  \item \texttt{CgZ5} = $(i\gamma_2 \gamma_8) \cdot \gamma_{11}$
    $\;\leftrightarrow\;$ $C\gamma_z\gamma_5$.
\end{itemize}

The choice of interpolator enters the two-point contraction as the
$(\Cmat\Gammaint)_{ij}$ matrix applied to the propagator indices $i,j$ and $k,l$.

% ======================================================================
\section{Connected Three-Point qTMD Function}
\label{sec:three_point}
% ======================================================================

\subsection{Definition}

The connected three-point qTMD function for a flavor insertion $f \in \{U, D\}$ is:

\begin{align}
  C_3^{\Gamma, f}(q, b, \tau; p_f, t_{\text{sep}}) =
    \sum_{\xvec} e^{-i q\cdot (\xvec - \xzerovec)} \;
    \operatorname{Tr}_{\text{spin,color}}\Bigl[
      &\mathrm{Seq}_f\bigl(\xvec; p_f, t_{\text{sep}}, \mathcal{P}\bigr) \notag \\
      &\times\; \Gamma \;
      \times\; \mathcal{O}_b\, S_q(\xvec, \xzerovec)
    \Bigr],
  \label{eq:C3_def}
\end{align}

where:
\begin{itemize}
  \item $\mathrm{Seq}_f(\xvec; p_f, t_{\text{sep}}, \mathcal{P})$ is the
    fixed-sink sequential propagator for flavor $f$, which already contains the
    baryon sink contraction, the final momentum projection $p_f$, the sink time
    $t_{\text{sep}}$, and the spin projection $\mathcal{P}$,
  \item $\Gamma$ is the bilinear insertion operator (one of the 16 gammas),
  \item $\mathcal{O}_b$ is the nonlocal displacement operator shifting the forward
    propagator by impact parameter $b = (b_\perp, b_z)$,
  \item $q = p_f - p_i$ is the momentum transfer,
  \item $\tau$ labels the operator insertion time slice.
\end{itemize}

\subsection{Data Layout}

The sequential propagator after the final $\gamma_5$ conjugation has shape:
\begin{equation}
  [\text{polarization},\, \omega,\, t,\, z,\, y,\, x_{cb},\,
    \mathrm{spin}_a,\, \mathrm{spin}_b,\, \mathrm{color}_a,\, \mathrm{color}_b].
\end{equation}

The forward propagator (with the nonlocal shift) has shape:
\begin{equation}
  [\omega,\, t,\, z,\, y,\, x_{cb},\,
    \mathrm{spin}_a,\, \mathrm{spin}_b,\, \mathrm{color}_a,\, \mathrm{color}_b].
\end{equation}

The three-point contraction einsum in the application is:
\begin{equation}
  \texttt{Seq[:, j, i, c, f]} \;\times\;
  \texttt{Gamma[g, i, m]} \;\times\;
  \texttt{S\_shifted[:, m, j, f, c]}
  \;\to\; [\text{pol}, g, \omega, t, z, y, x_{cb}],
\end{equation}
where indices are: $p$=polarization, $g$=gamma, $i$=spin\_src, $m$=spin\_sink,
$j$=spin\_b, $c$=color\_a, $f$=color\_b. This is followed by contraction with the
momentum phase $\Phi_q$ to project onto momentum transfer $q$.

% ======================================================================
\section{Fixed-Sink Sequential Sources}
\label{sec:sequential}
% ======================================================================

The fixed-sink sequential sources for the proton are constructed by
\texttt{create\_bw\_seq\_pyquda} in \texttt{bw\_seq\_pyquda.py}. The construction
follows these steps for each polarization $\mathcal{P}$ and flavor:

\begin{enumerate}
  \item \textbf{Boosted smearing} at the sink: $\texttt{prop} \to
    \texttt{boosted\_smearing}(\texttt{prop}, w, \texttt{boost\_out})$.
  \item \textbf{Baryon contraction}: depending on the flavor, either
    \texttt{up\_quark\_insertion\_pyquda} or \texttt{down\_quark\_insertion\_pyquda}
    contracts the propagator with epsilon tensors and the polarization matrix
    $\mathcal{P}$ to form the sequential source at the sink time slice.
  \item \textbf{Time slicing}: the source is restricted to time slice
    $t_{\text{sink}} = (t_{\text{src}} + t_{\text{sep}}) \bmod L_t$ using
    \texttt{sequential12}.
  \item \textbf{Momentum phase and $\gamma_5$}: the source is multiplied by
    $e^{i p_f \cdot (\xvec - \xzerovec)}$ and left-multiplied by $\gamma_5$:
    \begin{equation}
      \texttt{seq\_data} = \gamma_5 \cdot \Phi_{p_f} \cdot
        \texttt{seq\_data}^\dagger.
    \end{equation}
  \item \textbf{Inversion}: the source is smeared (boosted) and inverted to obtain
    the sequential propagator.
  \item \textbf{Post-inversion $\gamma_5$}: the sequential propagator is conjugated
    and multiplied by $\gamma_5$ on the right:
    \begin{equation}
      \texttt{final} = \texttt{prop\_inv}^\dagger \cdot \gamma_5.
    \end{equation}
\end{enumerate}

\subsection{Up-Quark Insertion (Flavor 1)}

The up-quark insertion function \texttt{up\_quark\_insertion\_pyquda} constructs the
sequential source from two up-quark propagators $Q_u$, one down-quark propagator
$Q_d$, the interpolator $\Gamma$, and the spin projection $\mathcal{P}$.

Let $Q_u$ and $Q_d$ be reshaped to $[\mathrm{vol}, 4, 4, 3, 3]$ with indices
$[\text{site},\, \text{spin}_{\text{sink}},\, \text{spin}_{\text{src}},\,
  \text{color}_{\text{sink}},\, \text{color}_{\text{src}}]$.

Internally, the following intermediate combinations are precomputed:
\begin{align}
  G^T D G &: \quad (\Gamma^T)_{ij}\, (Q_d)_{jkab}\, (\Gamma)_{kl}
    \;\to\; [\mathrm{vol}, i, l, a, b], \\
  \mathcal{P} D_u &: \quad (\mathcal{P})_{ij}\, (Q_u)_{jkab}
    \;\to\; [\mathrm{vol}, i, k, a, b], \\
  D_u \mathcal{P} &: \quad (Q_u)_{jkab}\, (\mathcal{P})_{kl}
    \;\to\; [\mathrm{vol}, j, l, a, b], \\
  \operatorname{Tr}[D_u \mathcal{P}] &: \quad
    (Q_u)_{kjab}\, (\mathcal{P})_{jk}
    \;\to\; [\mathrm{vol}, a, b].
\end{align}

The four epsilon-contracted terms are:
\begin{align}
  R_1 &: \quad
    \eps_{abc}\,\eps_{def}\,
    (G^T D G)_{mnbe}\, (Q_u)_{mnad}
    \;\to\; [\mathrm{vol}, b, e, a, d] \;
    \xrightarrow{\times \eps\eps} \; [\mathrm{vol}, c, f] \;
    \xrightarrow{\times \mathcal{P}} \; [\mathrm{vol}, i, j, c, f], \\[4pt]
  R_2 &: \quad
    \eps_{abc}\,\eps_{def}\,
    (\operatorname{Tr}[D_u\mathcal{P}])_{ad}\,
    (G^T D G)_{jibe}
    \;\to\; [\mathrm{vol}, i, j, c, f], \\[4pt]
  R_3 &: \quad
    \eps_{abc}\,\eps_{def}\,
    (\mathcal{P} D_u)_{ikad}\, (G^T D G)_{jkbe}
    \;\to\; [\mathrm{vol}, i, j, c, f], \\[4pt]
  R_4 &: \quad
    \eps_{abc}\,\eps_{def}\,
    (G^T D G)_{kiad}\, (D_u\mathcal{P})_{klbe}
    \;\to\; [\mathrm{vol}, i, l, c, f] \;
    \xrightarrow{i \leftrightarrow l} \; [\mathrm{vol}, i, j, c, f].
\end{align}

The total sequential source is:
\begin{equation}
  D_{\text{total}} = -(R_1 + R_2 + R_3 + R_4),
\end{equation}
with a final spin swap to reorder the output as
$[\mathrm{vol}, \mathrm{spin}_{\text{src}}, \mathrm{spin}_{\text{sink}},
  \mathrm{color}_{\text{sink}}, \mathrm{color}_{\text{src}}]$.

\subsection{Down-Quark Insertion (Flavor 2)}

The down-quark insertion \texttt{down\_quark\_insertion\_pyquda} is simpler (two
terms), since the down quark appears only once in the interpolating field. It takes a
single propagator $Q$ (isospin-symmetric), $\Gamma$, and $\mathcal{P}$.

The intermediate combinations are:
\begin{align}
  \mathcal{P} D &: \quad (\mathcal{P})_{ij}\, Q_{jiab}
    \;\to\; [\mathrm{vol}, a, b], \\
  G^T D G &: \quad (\Gamma^T)_{ij}\, Q_{jkab}\, (\Gamma)_{kl}
    \;\to\; [\mathrm{vol}, i, l, a, b], \\
  G^T D &: \quad (\Gamma^T)_{ij}\, Q_{jkab}
    \;\to\; [\mathrm{vol}, i, k, a, b], \\
  \mathcal{P} D G &: \quad (\mathcal{P})_{ij}\, Q_{jkab}\, (\Gamma)_{kl}
    \;\to\; [\mathrm{vol}, i, l, a, b].
\end{align}

The two terms are:
\begin{align}
  T_1 &: \quad
    \eps_{abc}\,\eps_{def}\,
    (\mathcal{P} D)_{fc}\, (G^T D G)_{uveb}
    \;\to\; [\mathrm{vol}, u, v, a, d], \\[4pt]
  T_2 &: \quad
    \eps_{abc}\,\eps_{def}\,
    (G^T D)_{ujec}\, (\mathcal{P} D G)_{jkfb}
    \;\to\; [\mathrm{vol}, u, k, a, d].
\end{align}

The total is:
\begin{equation}
  D_{\text{total}} = T_2 - T_1,
\end{equation}
with the same final spin swap.

\subsection{Spin Projections}

The available spin projection matrices $\mathcal{P}$ are defined as:

\begin{align}
  P_+ &= \frac{1}{4}(\gamma_0 + I), \\
  S_z^+ &= \gamma_0 - i\gamma_1\gamma_2, \qquad
  S_z^- = \gamma_0 + i\gamma_1\gamma_2, \\
  S_x^+ &= \gamma_0 - i\gamma_2\gamma_4, \qquad
  S_x^- = \gamma_0 + i\gamma_2\gamma_4.
\end{align}

The polarization projections stored in \texttt{PolProjections} are:
\begin{align}
  \texttt{PpUnpol} &= P_+, \\
  \texttt{PpSzp}   &= P_+ S_z^+, \\
  \texttt{PpSzm}   &= P_+ S_z^-, \\
  \texttt{PpSxp}   &= P_+ S_x^+, \\
  \texttt{PpSxm}   &= P_+ S_x^-.
\end{align}

% ======================================================================
\section{CG qTMD Operator}
\label{sec:cg_qtmd}
% ======================================================================

The Coulomb-gauge (CG) qTMD operator is a simple coordinate shift of the forward
propagator, without explicit gauge links:

\begin{equation}
  \mathcal{O}_b\, S_q(x, x_0) =
    S_q\bigl(x + b_\perp \cdot \ehat{\perp} + b_z \cdot \ehat{z},\; x_0\bigr),
  \label{eq:cg_shift}
\end{equation}

where $\ehat{\perp} \in \{\ehat{x}, \ehat{y}\}$ selects the transverse direction
(direction 0 or 1) and $\ehat{z}$ is the longitudinal direction (direction 2 on the
lattice). The displacement is parameterized by:
\begin{itemize}
  \item $b_\perp \in \{0, 1, \dots, b_T\}$ -- transverse displacement,
  \item $b_z \in \{-b_z, \dots, 0, \dots, b_z\}$ -- longitudinal displacement.
\end{itemize}

\subsection{Incremental Shift Strategy}

Rather than shifting the propagator by the full vector $(b_\perp, b_z)$ from the
origin for each Wilson-line index, the code maintains a running shifted propagator
and applies only incremental shifts:

\begin{equation}
  \texttt{prop\_shifted} =
    \texttt{prop\_prev.shift}(\Delta b_\perp, \text{dir}_\perp)\,
    \texttt{.shift}(\Delta b_z, \text{dir}_z),
\end{equation}

where $\Delta b_\perp = b_\perp^{\text{current}} - b_\perp^{\text{previous}}$ and
$\Delta b_z = b_z^{\text{current}} - b_z^{\text{previous}}$. This keeps successive
shifts small, minimizing communication overhead in the distributed lattice
implementation.

The method \texttt{create\_fw\_prop\_TMD\_CG} implements this logic: it receives the
current forward propagator, the current Wilson-line index $W$, and the previous
index, then applies the needed relative shifts.

% ======================================================================
\section{PDF Operators}
\label{sec:pdf}
% ======================================================================

The PDF operator is the $b_\perp = 0$ special case of the qTMD operator, involving
only longitudinal displacements along the $z$-direction.

\subsection{CG PDF}

The Coulomb-gauge PDF uses ordinary lattice shifts with $b_\perp = 0$:
\begin{equation}
  \mathcal{O}_{b_z}^{\text{CG}}\, S_q(x, x_0) =
    S_q(x + b_z \cdot \ehat{z},\; x_0).
\end{equation}
This is implemented via the same \texttt{create\_fw\_prop\_TMD\_CG} method with
$b_\perp = 0$ and $b_\perp^{\text{previous}} = 0$.

\subsection{GI PDF}

The gauge-invariant (GI) PDF uses covariant shifts along the $z$-direction, building
an explicit straight Wilson line through successive gauge-covariant nearest-neighbor
shifts:

\begin{equation}
  \mathcal{O}_{b_z}^{\text{GI}}\, S_q(x, x_0) =
    U_z(x + (b_z-1)\ehat{z}) \cdots U_z(x + \ehat{z})\, U_z(x)\;
    S_q(x + b_z \cdot \ehat{z},\; x_0).
\end{equation}

The implementation in \texttt{create\_fw\_prop\_PDF\_GI} uses
\texttt{gauge.pure\_gauge.covDev} to apply $\psi'(x) = U_\mu(x)\, \psi(x + \hat{\mu})$
for each nearest-neighbor step:
\begin{itemize}
  \item direction 2: $+z$ (forward covariant derivative),
  \item direction 6: $-z$ (backward covariant derivative).
\end{itemize}

The code enforces nearest-neighbor steps only ($\Delta b_z \in \{-1, 0, 1\}$) and
raises an error for larger jumps, protecting against accidentally skipping gauge
links. The shift is applied spinor-by-spinor and color-by-color in a loop over the
fermion field.

\subsection{Comparison}

\begin{center}
  \begin{tabular}{l|c|c|c}
    Variant & Gauge links & Transverse & Code method \\
    \hline
    CG\_TMD & None (CG) & $x$ and $y$ & \texttt{create\_fw\_prop\_TMD\_CG} \\
    CG\_PDF & None (CG) & none & \texttt{create\_fw\_prop\_TMD\_CG} \\
    GI\_PDF & Covariant & none & \texttt{create\_fw\_prop\_PDF\_GI}
  \end{tabular}
\end{center}

% ======================================================================
\section{Wilson Line Index Construction}
\label{sec:wilson}
% ======================================================================

\subsection{CG qTMD Wilson Indices}

The CG qTMD requires two sets of Wilson-line indices, one for each transverse direction:

\begin{itemize}
  \item \texttt{index\_list\_trans0}: transverse direction $0$ ($x$),
  \item \texttt{index\_list\_trans1}: transverse direction $1$ ($y$).
\end{itemize}

Each Wilson-line index is a 4-tuple $[b_T, b_z, \eta, \text{dir}]$ where:
\begin{itemize}
  \item $b_T \in [0, b_T^{\max}]$ -- transverse displacement,
  \item $b_z$ -- longitudinal displacement,
  \item $\eta = 0$ (irrelevant for CG, kept for interface compatibility),
  \item $\text{dir} \in \{0, 1\}$ -- transverse direction ($x$ or $y$).
\end{itemize}

The generation algorithm:
\begin{enumerate}
  \item For each $b_z \in [0, b_z^{\max}]$, $b_T \in [0, b_T^{\max}]$:
    \begin{itemize}
      \item add $[b_T, b_z, 0, \text{dir}]$,
      \item if $b_z \neq 0$, also add $[b_T, -b_z, 0, \text{dir}]$.
    \end{itemize}
  \item Reorder the index list to minimize adjacent differences in $(b_T, b_z)$.
    The reordering sorts by $(b_T, b_z)$ first, then locally permutes to make each
    step as small as possible.
\end{enumerate}

The final combined list is:
\begin{equation}
  \texttt{W\_index\_list\_CG} =
    \texttt{index\_list\_trans0} \;+\; \texttt{index\_list\_trans1}.
\end{equation}

\subsection{PDF Wilson Indices}

The PDF Wilson-line index list contains only $b_T = 0$ entries:
\begin{enumerate}
  \item For each $b_z \in [0, b_z^{\max}]$: add $[0, b_z, 0, 0]$.
  \item For each $b_z \in [1, b_z^{\max}]$: add $[0, -b_z, 0, 0]$.
\end{enumerate}

Since $b_T = 0$ for all entries, the transverse direction is always 0 (dummy).
The positive and negative $b_z$ values are separated to facilitate the incremental
shift strategy: for GI\_PDF, the gauge-covariant link construction requires resetting
to the unshifted propagator when the sign of $b_z$ changes (i.e., when $b_z = 0$ or
$b_z = -1$ in the loop).

% ======================================================================
\section{Code Implementation}
\label{sec:implementation}
% ======================================================================

The measurement is organized into a class \texttt{proton\_TMD} in
\texttt{proton\_qTMD\_pyquda.py}.

\subsection{Constructor Parameters}

\begin{center}
  \begin{tabular}{l|l}
    Parameter & Description \\
    \hline
    \texttt{eta} & List of $\eta$ values (irrelevant for CG) \\
    \texttt{b\_z} & Maximum longitudinal displacement \\
    \texttt{b\_T} & Maximum transverse displacement \\
    \texttt{pilist} & Initial momenta $p_i$ for 2pt \\
    \texttt{width} & Gaussian smearing width \\
    \texttt{boost\_out} & C2 sink smearing \\
  \end{tabular}
\end{center}

The application uses \texttt{boost\_in} to smear the source before the forward
inversion.  The production measurement object stores \texttt{boost\_out} and
uses it for the sink side of the two-point contraction; the same
\texttt{boost\_out} is passed to the fixed-sink sequential-source
construction.  Thus the saved two-point function has the endpoint convention
\begin{equation}
  C_{2}^{\mathrm{out,in}}(p,t)
  =
  \left\langle
    \chi_{S_{\mathrm{out}}}(p,t)\,
    \overline{\chi}_{S_{\mathrm{in}}}(0)
  \right\rangle .
\end{equation}
The current canonical applications set
\(\texttt{boost\_in}=\texttt{boost\_out}=(0,0,0)\), so this correction does
not change their standard zero-boost data.  Manually produced older
unequal-boost two-point data used \texttt{boost\_in} at both endpoints and
must be recomputed.

\subsection{Method Summary}

\begin{itemize}
  \item \texttt{contract\_2pt\_TMD(latt\_info, prop\_f, phases, tag, interpolator)}:
    Full two-point contraction with all 16 gammas (Section~\ref{sec:contraction_code}).
  \item \texttt{create\_TMD\_Wilsonline\_index\_list\_CG()}:
    Generate CG Wilson indices for both transverse directions (Section~\ref{sec:wilson}).
  \item \texttt{create\_PDF\_Wilsonline\_index\_list()}:
    Generate PDF Wilson indices ($b_\perp = 0$, Section~\ref{sec:wilson}).
  \item \texttt{create\_fw\_prop\_TMD\_CG(prop\_f, W\_index, WL\_indices\_previous)}:
    Incremental CG shift (Section~\ref{sec:cg_qtmd}).
  \item \texttt{create\_fw\_prop\_PDF\_GI(gauge, prop\_f, W\_index, WL\_indices\_previous)}:
    Gauge-covariant GI shift (Section~\ref{sec:pdf}).
\end{itemize}

% ======================================================================
\section{Array and Output Shape Convention}
\label{sec:output}
% ======================================================================

\subsection{Internal Array Shapes}

\begin{center}
  \begin{tabular}{l|l}
    Object & Shape \\
    \hline
    Propagator data & $[\omega, t, z, y, x_{cb}, s_a, s_b, c_a, c_b]$ \\
    Momentum phases & $[p, \omega, t, z, y, x]$ or $[p, \omega, t, z, y]$ \\
    Gamma array & $[g, m, n]$ (16 $\times$ 4 $\times$ 4) \\
    Sequential propagator & $[\text{pol}, \omega, t, z, y, x_{cb}, s_a, s_b, c_a, c_b]$ \\
    C2 result & $[g, p, t]$ \\
    C3 result & $[\text{WL}, \text{pol}, q, g, t]$
  \end{tabular}
\end{center}

where:
\begin{itemize}
  \item $\omega$: even-odd parity ($L_t \times L_z \times L_y \times L_x/2$),
  \item $x_{cb}$: checkerboarded $x$ index,
  \item $s_a, s_b$: spin indices (0--3),
  \item $c_a, c_b$: color indices (0--2),
  \item $g$: gamma matrix index (0--15),
  \item $p, q$: momentum indices,
  \item $\text{pol}$: polarization index,
  \item $\text{WL}$: Wilson-line index.
\end{itemize}

\subsection{HDF5 Output Layout}

\subsubsection{Two-Point Function}

\begin{verbatim}
  <tag>.h5
    /SS/<gamma>/PX..PY..PZ..  ->  dataset of shape [Lt]
\end{verbatim}

Each gamma channel stores all momenta as separate datasets. The time axis is rolled
so that $t=0$ corresponds to the source position.

\subsubsection{Three-Point qTMD/PDF}

\begin{verbatim}
  <tag>.h5
    /SS/<gamma>/PX..PY..PZ../
      <b_X or b_Y>/eta0/bT<bT>/bz<bz>  ->  dataset of shape [tsep+2]
\end{verbatim}

The Wilson-line index $[b_T, b_z, \eta, \text{dir}]$ maps to the HDF5 group path:
\begin{itemize}
  \item \texttt{dir} = 0 $\to$ \texttt{b\_X}, \texttt{dir} = 1 $\to$ \texttt{b\_Y},
  \item \texttt{eta} = 0 $\to$ \texttt{eta0},
  \item $b_T \to$ \texttt{bT}$\langle b_T\rangle$,
  \item $b_z \to$ \texttt{bz}$\langle b_z\rangle$.
\end{itemize}

For PDF data, only \texttt{b\_X} appears (since $b_T = 0$ always, and the transverse
direction is dummy).

The flavor-specific tags encode the flavor as part of the AMA field:
\texttt{"CG.D.ex"} for down-quark CG insertion, \texttt{"CG.U.ex"} for up-quark
CG insertion.

% ======================================================================
\section{Key Differences from Pion qTMD}
\label{sec:differences}
% ======================================================================

The proton qTMD measurement differs from the pion qTMD in several fundamental ways:

\begin{center}
  \begin{tabular}{l|l|l}
    Aspect & Pion qTMD & Proton qTMD \\
    \hline
    Contraction & $\gamma_5$-hermiticity & Baryon epsilon + 2 exchange terms \\
    Quark lines & 2 (one contraction term) & 3 (two contraction terms) \\
    Flavor dim. & None (single channel) & $[U, D]$ via 2 seq.\ sources \\
    Polarization & None & Via spin projectors $\mathcal{P}$ \\
    Propagators & Separate fw/bw & Same prop reused (isospin) \\
    Mom.\ params. & \texttt{pos\_boost}/\texttt{neg\_boost} &
      \texttt{boost\_in}/\texttt{boost\_out} \\
    Seq.\ source & 1 meson seq.\ source & 2 baryon seq.\ sources \\
    Output shape & $[\text{WL}, q, g, t]$ &
      $[\text{WL}, \text{pol}, q, g, t]$
  \end{tabular}
\end{center}

The isospin-symmetric approximation ($m_u = m_d$) means the same forward propagator
is used for all three valence quark lines in the two-point function and for the
sequential source construction. The up and down sequential sources differ only in
their baryon contraction pattern, not in the input propagators.

% ======================================================================
\section{Limitations}
\label{sec:limitations}
% ======================================================================

The current implementation has the following limitations:

\begin{enumerate}
  \item \textbf{Connected diagrams only.} Disconnected insertions (where the operator
    couples to a sea-quark loop) are not included. These contribute to the singlet
    TMDs and may be substantial at low $x$.
  \item \textbf{Isospin symmetry.} The code assumes $m_u = m_d$ and reuses the same
    forward propagator for all quark lines. Isospin-breaking effects from
    $m_u \neq m_d$ or from electromagnetic corrections are not captured.
  \item \textbf{No flavor mixing or renormalization.} The bare lattice operators are
    not matched to the $\overline{\text{MS}}$ scheme. Renormalization factors and
    operator mixing (especially for the Wilson-line operators) must be applied
    externally.
  \item \textbf{CG paths have no explicit gauge links.} The Coulomb-gauge qTMD
    operator is a simple coordinate shift. This is gauge-dependent unless the
    lattice is fixed to Coulomb gauge. Residual gauge-fixing uncertainties are
    not quantified.
  \item \textbf{No disconnected TMD contributions.} The full proton TMD includes
    disconnected diagrams (operator insertion on a sea-quark loop connecting to
    the proton via gluons), which are not computed.
  \item \textbf{Fixed sink only.} The sequential-source method fixes the sink time
    and momentum. To vary $t_{\text{sep}}$ or $p_f$, separate inversions are required.
\end{enumerate}

% ======================================================================
\section{Connected GI qTMD Update}
\label{sec:connected_gi_qtmd}
% ======================================================================

The connected proton qTMD code has a gauge-invariant staple path in addition
to the Coulomb-gauge coordinate-shift qTMD and straight PDF operators.  Both
Perlmutter and Aurora call the single production implementation in
\texttt{application/nucleon\_TMD/shared\_runner.py}; platform entry points only
select backend and ensemble defaults.

The connected \texttt{GI\_qTMD} operator uses the same Wilson index convention as
the pion and disconnected qTMD code,
\begin{equation}
  W = [b_T,\; b_z,\; \eta,\; d_\perp],
\end{equation}
with $d_\perp=0,1$ for the transverse $x$ or $y$ direction.  The fixed-length
staple path is
\begin{align}
  x
  &\rightarrow x + \left(\eta + \frac{b_z}{2}\right)\hat z, \\
  &\rightarrow x + \left(\eta + \frac{b_z}{2}\right)\hat z
       + b_T \hat e_{d_\perp}, \\
  &\rightarrow x + b_z \hat z + b_T \hat e_{d_\perp}.
\end{align}
The total staple length is therefore
\begin{equation}
  L_{\rm staple} = 2\eta + b_T,
\end{equation}
which is independent of $b_z$.  The allowed connected \texttt{GI\_qTMD}
indices satisfy
\begin{equation}
  b_z \in 2\mathbb{Z}, \qquad \eta \geq \frac{|b_z|}{2}, \qquad b_T \geq 0.
\end{equation}

The source module exposes this through two new helpers:
\begin{itemize}
  \item \texttt{create\_TMD\_Wilsonline\_index\_list\_GI()} generates the
    connected staple-index list.
  \item \texttt{create\_fw\_prop\_TMD\_GI(gauge, prop, W, staple\_links)}
    applies a cached gauge-only staple transporter to each spin-color source
    column of the forward propagator.
\end{itemize}
The actual baryon contraction is unchanged: the shifted or covariantly
transported forward line is contracted with the up- and down-insertion
sequential propagators,
\begin{equation}
  C_{3,g}^{f}(q,W,\tau)
  =
  \sum_{\mathbf{x}}
  e^{-i\mathbf{q}\cdot(\mathbf{x}-\mathbf{x}_0)}
  \mathrm{Tr}\!\left[
    \mathrm{Seq}_{f}(\mathbf{x},\tau;p_f,t_{\rm sep},\mathcal{P})\,
    \Gamma_g\,
    \mathcal{O}_{W}^{\rm GI} S(\mathbf{x},\tau;\mathbf{x}_0,0)
  \right],
\end{equation}
where $f=U,D$ labels the connected flavor insertion and $\mathcal{P}$ is the
spin projector.

Output files use the tags
\begin{equation}
  \texttt{GI\_qTMD.U.ex}, \qquad \texttt{GI\_qTMD.D.ex},
\end{equation}
while the original CG files keep the tags \texttt{CG.U.ex} and
\texttt{CG.D.ex}.  The HDF5 internal path still follows the common qTMD writer
layout based on $[b_T,b_z,\eta,d_\perp]$.

% ======================================================================
\section{Application Workflow Summary}
\label{sec:workflow}
% ======================================================================

The production applications follow this workflow for each source position:

\begin{enumerate}
  \item \textbf{Forward propagator}: Generate a point source at position
    $(x_0, y_0, z_0, t_0)$, apply boosted smearing with
    \texttt{boost\_in}, and invert.
  \item \textbf{Two-point contraction}: Call \texttt{contract\_2pt\_TMD} with the
    forward propagator and momentum phases, applying sink smearing with
    \texttt{boost\_out}. This computes all 16 gamma channels and saves the
    result as HDF5.
  \item \textbf{Sequential sources}: Build two sequential backward propagators
    via \texttt{create\_bw\_seq\_pyquda}, using \texttt{boost\_out} for the
    fixed sink:
    \begin{itemize}
      \item Flavor 2 (down): \texttt{sequential\_bw\_prop\_down},
      \item Flavor 1 (up): \texttt{sequential\_bw\_prop\_up}.
    \end{itemize}
  \item \textbf{TMD contractions (CG)}: Loop over Wilson-line indices for
    directions $+x$ and $+y$, incrementally shifting the forward propagator
    and contracting with both sequential propagators and the 16 gammas.
    Save results per polarization, gamma, and flavor.
  \item \textbf{TMD contractions (GI)}: In the new
    \texttt{application/nucleon\_TMD} workflows, build cached fixed-length
    staple transporters, apply them to the forward propagator, and contract
    with the same sequential propagators.  This path writes
    \texttt{GI\_qTMD.U.ex} and \texttt{GI\_qTMD.D.ex}.
  \item \textbf{PDF contractions (GI)}: Loop over PDF Wilson-line indices,
    apply gauge-covariant $z$-shifts via \texttt{create\_fw\_prop\_PDF\_GI},
    and contract similarly.
  \item \textbf{PDF contractions (CG)}: Same as GI but using ordinary lattice
    shifts.
\end{enumerate}

The momentum transfer $q$ for the three-point functions is constructed from the
\texttt{qext} parameter list: $q = (q_x, q_y, q_z, 0)$ with $q_z = 0$ for TMD
and unrestricted for PDF. The initial momentum is $p_i = p_f - q$, and the
two-point momenta are the negative of the \texttt{p\_2pt} list.

On Perlmutter, the smoke-test entry point is
\begin{equation}
  \texttt{run\_nucleon\_TMD.sh --config\_num CFG
  --mg-block 8.8.4.4}.
\end{equation}
\texttt{--config\_num} is mandatory and is never read from the environment.
Unknown arguments fail before gauge I/O.  The MG block is an explicit runtime
choice; multiple levels use semicolon separation and \texttt{none} disables
MG.  MG, solver tolerance, and maximum iterations are not HDF5 provenance and
do not enter the sample-log identity.

At startup rank zero reads the sample log using exact-line matching and
broadcasts the completed-source set.  Completed sources are skipped before
source construction and inversion.  A source is appended only after C2 and
all products enabled in that invocation close successfully.  The log is the
only resume state and is intentionally independent of HDF5 location.  It may
therefore be reused only when operator switches, momentum/Wilson grids,
smearing, interpolator, and sink separation are unchanged.

The main environment toggles are
\texttt{NUCLEON\_TMD\_RUN\_CG\_QTMD},
\texttt{NUCLEON\_TMD\_RUN\_GI\_QTMD}, and
\texttt{NUCLEON\_TMD\_RUN\_PDF}.  The staple parameters are controlled by
\texttt{NUCLEON\_TMD\_ETA}, \texttt{NUCLEON\_TMD\_BZ}, and
\texttt{NUCLEON\_TMD\_BT}.  For the fixed-length GI qTMD path, use even
\texttt{NUCLEON\_TMD\_BZ} and choose
\texttt{NUCLEON\_TMD\_ETA >= abs(NUCLEON\_TMD\_BZ)/2}.

% ======================================================================
\section*{Acknowledgments}
% ======================================================================

This code was developed as part of the PyQUDA measurement framework for lattice QCD
nucleon structure calculations. The contraction strategy and sequential-source
construction build on established methods in the baryon spectroscopy literature.

\end{document}
